\documentclass[12pt]{article}
\usepackage[T1]{fontenc}
\usepackage[utf8]{inputenc}
\usepackage[brazil]{babel}
\usepackage[margin=1in]{geometry}
\usepackage{ragged2e}
\usepackage[normalem]{ulem}
\setlength{\parindent}{0pt}
\setlength{\parskip}{0.8\baselineskip}
\begin{document}
\justifying
\noindent Verifica-se que a questão objeto da controvérsia jurídica suscitada no recurso manejado pela parte recorrente esteve afetada no representativo da controvérsia - \textbf{Tema n. }\textbf{358} da TNU - (PU no processo n. PEDILEF 0500179-22.2022.4.05.8311/PE),  afetado em 13/03/2024 foi julgado (trânsito em julgado em 29/11/2024), por meio do qual foi elucidada a seguinte questão:

\noindent \textit{"Saber se, para fins de concessão de aposentadoria por idade urbana com DER após a EC 103/2019, permanece a necessidade de cumprimento do requisito da carência, particularmente para quem precisa usar a regra de transição do art. 18 da EC 103, ou se a regra de transição prevista no art. 18, da EC 103/19 não exige mais tal requisito (bastando ao beneficiário preencher, cumulativamente, os requisitos "idade" e "tempo de contribuição"), de forma que as contribuições recolhidas em atraso pelo contribuinte individual possam ser computados como tempo de contribuição (ainda que este tenha perdido a qualidade de segurado)."}

\noindent Restou firmada a seguinte tese:

\noindent \textit{"1. Tempo de contribuição e carência são institutos distintos. 2. Carência condiz com contribuições tempestivas. 3. O art. 18 da EC 103/2019 não dispensa a carência para a concessão de aposentadoria." (Tema 358 TNU).}

\noindent Com efeito, depreende-se da leitura do voto condutor do acórdão que o mesmo se encontra em conformidade com o decidido no  tema pela TNU.

\noindent A proposito, consta do acordao hostilizado: (...) Aposentadoria por Tempo de Contribuição (Art. 55/6), Benefícios em Espécie, DIREITO PREVIDENCIÁRIO

\noindent Nessas condições, o conteúdo do Acórdão recorrido é consentâneo com o entendimento da TNU, firmado em julgamento de recurso representativo de controvérsia, o que atrai a incidência da regra prevista pelo art. 14, III, \textit{b,} da Resolução CJF n. 586/2019:

\noindent \textit{\textbf{Art. 14.}}\textit{ Decorrido o prazo para contrarrazões, os autos serão conclusos ao magistrado responsável pelo exame preliminar de admissibilidade, que deverá, de forma sucessiva: (...) }\textit{\textbf{III }}\textit{- negar seguimento a pedido de uniformização de interpretação de lei federal interposto contra acórdão que esteja em conformidade com entendimento consolidado: }\textit{\textbf{b)}}\textit{ em }\uline{\textit{\textbf{recurso representativo de controvérsia pela Turma Nacional de Uniformização}}}\textit{ ou em pedido de uniformização de interpretação de lei dirigido ao Superior Tribunal de Justiça;(...).}

\noindent Posto isso, com arrimo no \uline{\textbf{art. 14, III, }}\uline{\textit{\textbf{b}}}\textit{,} da Resolução CJF n. 586/2019 c/c art. 1.030, I, do CPC, \textbf{NEGO SEGUIMENTO }\textbf{a}o\textbf{ PU}\textbf{.}

\noindent Intimem-se. Aguarde-se o decurso do prazo recursal. Após, certifique-se o trânsito em julgado e devolva-se ao juízo de origem.
\end{document}
